\section{Qualifikationen}
\cvitem{05/2017}{Promotion: „Philosophie und Philosophieren im Unterricht. Empirische Erschließung einer widersprüchlichen Praxis.“ Goethe-Universität Frankfurt am Main.}
\cvitem{04/2010}{Staatsexamen für das Lehramt an Gymnasien mit den Fächern: Politik und Wirtschaft, Philosophie und Ethik, Goethe-Universität Frankfurt am Main.}
\cvitem{10/2006}{Diplom in Pädagogik, Goethe-Universität Frankfurt am Main.}
\cvitem{06/1999}{Abitur, Goethe Gymnasium Neu-Isenburg.}

\section{Berufliche Laufbahn}
\cvitem{seit 04/2024}{Alpen-Adria-Universität Klagenfurt, Institut für Unterrichts- und Schulentwicklung – Senior Scientist}
\cvitem{04/2021–03/2024}{Goethe-Universität, Fachbereich Erziehungswissenschaften – Wissenschaftlicher Mitarbeiter}
\cvitem{10/2019–03/2021}{Goethe-Universität, Fachbereich Erziehungswissenschaften – Vertretung der W3-Professur „Sozialpädagogik und Familienforschung“}
\cvitem{10/2013–09/2019}{Goethe-Universität, Fachbereich Erziehungswissenschaften – Wissenschaftlicher Mitarbeiter}
\cvitem{2012–2013}{Goethe-Universität, Fachbereich Erziehungswissenschaften – Wissenschaftliche Hilfskraft}
\cvitem{2010–2012}{Studienseminar Offenbach, Stadthof 13, 63065 Offenbach – Studienreferendar für das Lehramt an Gymnasien}
\cvitem{2009–2010}{Goethe-Universität, Zentrum für Lehrerbildung – Wissenschaftliche Hilfskraft}
\cvitem{2002–2010}{Verein für Psychotherapie, Beratung und Heilpädagogik, Alexanderstr. 29, 60489 Frankfurt – Pädagogischer Mitarbeiter}
\cvitem{1999–2002}{Heilpädagogisches Institut Vincenzhaus, Schule für Erziehungshilfe und Kranke, Vincenzstr. 29, 65719 Hofheim – Zivildienst und Pädagogischer Mitarbeiter}

\section{Verzeichnis der Lehrveranstaltungen}
\cvsubhead{Wintersemester 2023/2024 (Ausblick)}
\begin{itemize}
\item Einführung in das Narrative Interview
\item Learning Democracy in Zeiten gesellschaftlicher Krisen III
\end{itemize}
\cvsubhead{Sommersemester 2023}
\begin{itemize}
\item Einführung in das Narrative Interview
\item Learning Democracy in Zeiten gesellschaftlicher Krisen II
\end{itemize}
\cvsubhead{Wintersemester 2022/2023}
\begin{itemize}
\item Befreit von der Lehrverpflichtung für Arbeiten in einem interdisziplinären Forschungsverbund (Robust Nature).
\end{itemize}
\cvsubhead{Sommersemester 2022}
\begin{itemize}
\item Learning Democracy in Zeiten gesellschaftlicher Krisen
\end{itemize}
\cvsubhead{Wintersemester 2021/2022}
\begin{itemize}
\item Einführung in erziehungswissenschaftliche Grundbegriffe
\item Bildung zu planetarem Denken als Antwort auf die Herausforderung des Anthropozäns?
\end{itemize}
\cvsubhead{Sommersemester 2021}
\begin{itemize}
\item Zur materialistischen Pädagogik II
\item Forschung im Bereich Education for Sustainability
\item Einführung in die Objektive Hermeneutik
\item Zum Verhältnis von Bildung in Müdigkeit und Erziehung zum gesellschaftlichen Zusammenhalt
\end{itemize}
\cvsubhead{Wintersemester 2020/2021}
\begin{itemize}
\item \textbf{Vorlesung:} Zur Geschichte, den Grundbegriffen und den Grundfragen der Erziehungswissenschaft
\item Zur materialistischen Pädagogik III
\item Einführung in die Dokumentarische Methode
\item Empowerment und nachhaltige Stadtentwicklung durch Forschung mittels der Dokumentarischen Methode
\end{itemize}
\cvsubhead{Sommersemester 2020}
\begin{itemize}
\item Zur materialistischen Pädagogik II
\item Forschung im Bereich Education for Sustainability
\item Einführung in die Objektive Hermeneutik
\item Zum Verhältnis von Bildung in Müdigkeit und Erziehung zum gesellschaftlichen Zusammenhalt
\end{itemize}
\cvsubhead{Wintersemester 2019/2020}
\begin{itemize}
\item Bildung für nachhaltige Entwicklung und Transformation der Schule
\item Einführung in die Erziehungswissenschaft
\item Ringvorlesung Bildung für nachhaltige Entwicklung
\item Zur materialistischen Pädagogik I
\item Bildung für nachhaltige Entwicklung – Zur Erforschung der Implementierung eines pädagogischen Programms
\end{itemize}
\cvsubhead{Sommersemester 2019}
\begin{itemize}
\item Education for Sustainability V
\item Analyse von Unterrichtsprozessen am Beispiel des Philosophie-Unterrichts III
\end{itemize}
\cvsubhead{Wintersemester 2018/2019}
\begin{itemize}
\item Zum Verhältnis von Erziehungstheorie und Politischer Philosophie
\item Education for Sustainability IV
\item Zur kritischen Bildungs- und Erziehungstheorie H.-J. Heydorns II
\item Zur pädagogischen Professionalität in der (außer-)schulischen Politischen Bildung und Umweltbildung
\item Schulpraktische Studien Nachbereitung (forschungsbezogen)
\end{itemize}
\cvsubhead{Sommersemester 2018}
\begin{itemize}
\item Analysen von Unterrichtsprozessen unter besonderer Berücksichtigung der Thematik „Bildung für nachhaltige Entwicklung“
\item Education for Sustainability III
\item Critical Pedagogy
\item Schulpraktische Studien Vorbereitung (forschungsbezogen)
\end{itemize}
\cvsubhead{Wintersemester 2017/2018}
\begin{itemize}
\item Analyse von Unterrichtsprozessen unter besonderer Berücksichtigung der Ontogenese Bürgerlicher Kälte
\item Education for Sustainability II
\item Critical Pedagogy
\item Lehrauftrag an der Europa-Universität Flensburg, Philosophisches Seminar: Praktikumsbegleitendes Seminar
\end{itemize}
\cvsubhead{Sommersemester 2017}
\begin{itemize}
\item John Dewey: Democracy and Education
\item Education for Sustainability
\item Negative Pädagogik
\item Analysen von Unterrichtsprozessen unter besonderer Berücksichtigung der Thematik „Bildung für nachhaltige Entwicklung“
\end{itemize}
\cvsubhead{Wintersemester 2016/2017}
\begin{itemize}
\item Analysen von Unterrichtsprozessen unter besonderer Berücksichtigung der Thematik „Bildung für nachhaltige Entwicklung“
\item Lektüreseminar H.-J. Heydorn „Über den Widerspruch von Bildung und Herrschaft“
\item Bildung für nachhaltige Entwicklung
\item Zur kritischen Bildungs- und Erziehungstheorie H.-J. Heydorns
\end{itemize}
\cvsubhead{Sommersemester 2016}
\begin{itemize}
\item Lehrerbildung mittels fallrekonstruktiver Unterrichtsforschung
\item Lektüreseminar Theodor W. Adorno
\end{itemize}
\cvsubhead{Wintersemester 2015/2016}
\begin{itemize}
\item Unterrichten: Analyse von Unterrichtsprozessen II (Erziehung, Bildung, Didaktik)
\item Lektüreseminar: Wilhelm von Humboldt
\end{itemize}
\cvsubhead{Sommersemester 2015}
\begin{itemize}
\item Unterrichten. Analyse von Unterrichtsprozessen (Erziehung, Bildung, Didaktik)
\item Lektüreseminar: Johann Friedrich Herbart
\end{itemize}
\cvsubhead{Wintersemester 2014/2015}
\begin{itemize}
\item Analyse von Unterrichtsprozessen am Beispiel des Philosophie-Unterrichts
\item Lektüreseminar: Die Pädagogik der Philanthropen
\end{itemize}
\cvsubhead{Sommersemester 2014}
\begin{itemize}
\item Analyse von Unterrichtsprozessen an Beispielen des Philosophie-Unterrichts II
\item Lektüreseminar: Johann Heinrich Pestalozzi
\end{itemize}
\cvsubhead{Wintersemester 2013/2014}
\begin{itemize}
\item Emil oder über die Erziehung
\item Analyse von Unterrichtsprozessen an Beispielen des Philosophie-Unterrichts I
\end{itemize}
\cvsubhead{Lehraufträge}
\cvitem{Herbstsem. 2017}{Praktikumsbegleitendes Seminar. Universität Flensburg, Philosophisches Seminar.}
\cvitem{Frühjahrssem. 2019}{Übergang in eine transformierte Gesellschaft in und durch Bildung für nachhaltige Entwicklung in Schule und Unterricht. Universität Flensburg, Institut für Erziehungswissenschaften.}
\cvitem{Sommersem. 2019}{Erziehung zur Mündigkeit. Universität Flensburg, Institut für Erziehungswissenschaften; Bildung für nachhaltige Entwicklung. Universität Klagenfurt (Österreich); Unterricht beobachten, rekonstruieren, initiieren. Universität Mainz.}

\cvausblick{Ausblick}
\begin{itemize}
\item On educational theory of professionalisation and education for sustainable development – Wintersemester 2023/2024, Universität Marburg.
\end{itemize}
\section{Publikationsverzeichnis}
\subsection{Monographien}
\begin{enumerate}
\item Kminek, H. (2018): Philosophie und Philosophieren im Unterricht. Empirische Erschließung einer widersprüchlichen Praxis. (Dissertation) Opladen, Berlin, Toronto.
\end{enumerate}
{\itshape\bfseries\small Monographien [in Vorbereitung]}\par\vspace{2pt}
\begin{enumerate}
\item Kminek, H. (2024): Zu den Aporien der Bildung für eine nachhaltige Entwicklung. (Arbeitstitel, in Vorbereitung)
\end{enumerate}
\subsection{Herausgeberschaften}
\cvsubhead{Herausgeberbände}
\begin{enumerate}
\item Kminek, H./Thein, C./Torkler, R. (Hrsg.) (2018): Zwischen Präskription und Deskription – zum Selbstverständnis der Philosophiedidaktik. Opladen, Berlin, Toronto: Budrich.
\item Kminek, H./Bank, F./Fuchs, L. (Hrsg.) (2020): Kontroverses Miteinander. Interdisziplinäre und kontroverse Positionen zur Bildung für eine nachhaltige Entwicklung. Norderstedt: BoD.
\item Kminek, H./Geyer, A./Siewert, M. (Hrsg.) (2022): Engaged Reflections – Reflected Engagement. Transdisciplinary Contributions to Education for Sustainable Development. Opladen, Berlin, Toronto: Budrich.
\item Kminek, H. (Hrsg.) (2024): Survival through Bildung. (Scientific Contribution to Didactics of Philosophy and Philosophy of Education, Bd. 8). Opladen, Farmington Hills: Verlag Barbara Budrich, 164 S.
\item Kminek, H./Fingerle, M./Andresen, S. (Hrsg.) (2026): Erziehungswissenschaft und Ethik. Verstrickungen, Widersprüche und Spannungsfelder in Forschung und Lehre. Stuttgart: W. Kohlhammer, 112 S.
\item Kminek, H./Singer-Brodowski, M./Holz, V. (Hrsg.) (2023): „Ökologische, gesellschaftliche und individuelle Umbrüche und ihre Bedeutung für Bildung für nachhaltige Entwicklung“. Opladen, Berlin, Toronto: Budrich. \textit{(in Vorbereitung)}
\item Ertl, H./Holz, V./Idel, T.-S./Kminek, H./Singer-Brodowski, M./Wulf, Ch. (Hrsg.) (2025): Bildung für nachhaltige Entwicklung. Einblicke in aktuelle Forschungsprojekte. (Edition Zeitschrift für Erziehungswissenschaft, Bd. 17), 259 S.
\item Costa, J./Kminek, H./Ruckelshauß, T./Singer-Brodowski, M./Weselek, J. (Hrsg.) (2025): Bildung für nachhaltige Entwicklung: Kontroversen und Debatten. (Schriftenreihe „Ökologie und Erziehungswissenschaft“ der Kommission Bildung für nachhaltige Entwicklung der DGfE). Opladen, Farmington Hills: Budrich, 166 S.
\item Holfelder, A.-K./Costa, J./Hufnagl, J./Kminek, H./Singer-Brodowski, M./Taube, D. (Hrsg.) (2026): Spannungsfelder im Forschungsfeld Bildung für nachhaltige Entwicklung. Ambivalenzen und Widersprüche in der Konzeption, Vermittlung und Aneignung von BNE. Opladen, Farmington Hills: Budrich, 138 S.
\end{enumerate}
\cvsubhead{Buchreihe}
\begin{enumerate}
\item Mitherausgeber der „Reihe – Wissenschaftliche Beiträge zur Philosophiedidaktik und Bildungsphilosophie“ beim Barbara Budrich-Verlag (gemeinsam mit Christian Thein, seit 2016).
\end{enumerate}
\cvsubhead{Zeitschrift}
\begin{enumerate}
\item Akbaba, Y./Kabel, S./Kminek, H./Martens, M./Meier, M./Pollmanns, M./Schindler, Ch. (Hrsg.): datum \& diskurs – Zeitschrift für Komplementär-, Sekundär- und Reanalysen im Feld der qualitativen Schul- und Unterrichtsforschung.
\end{enumerate}
\subsection{Aufsätze/Beiträge in Sammelwerken und Zeitschriften [double-blind peer-reviewed]}
\begin{enumerate}
\item Pollmanns, M./Leser, Ch./Kminek, H./Kabel, S./Hünig, R.: „Professionalisierung durch kasuistisch ausgerichtete Schulpraktische Studien? Analysen studentischer Fallarbeit.“ In: Fraefel, U./Seel, A. (Hrsg.) (2017): Schulpraktische Professionalisierung: Konzeptionelle Perspektiven. S. 179–194. Münster: Waxmann.
\item Veja, C./Sticht, K./Schindler, Ch./Kminek, H. (2017): SMW Based VRE for Addressing Multi-Layered Data Analysis – The Use Case of Classroom Interaction Interpretation. DOI: 10.1145/3125433.3125457.
\item Schindler, Ch./Kminek, H./Hocker, J./Meier, M./Veja, C. (2020): „Collaborative Open Analysis in a Qualitative Research Environment“. In: Education for Information. Engaging with Open Science in Teaching \& Learning, S. 1–15. DOI: 10.3233/EFI-190261.
\item Kminek, H./Meier, M./Schindler, Ch./Hocker, J./Veja, C. (2020): „Interpretieren im Kontext virtueller Forschungsumgebungen – zu den Potentialen und Grenzen des Einsatzes einer virtuellen Forschungsumgebung und ihren Einsatz in der akademischen Lehre“. In: Zeitschrift für Qualitative Forschung, 2/20 Methodenlernen.
\item Holfelder, A.-K./Singer-Brodowski, M./Holz, V./Kminek, H. (2020): Erziehungswissenschaftliche Fragen im Zusammenhang mit der Bewegung Fridays for Future. Zeitschrift für Pädagogik, 01/2021.
\item Kminek, H. (2022): About Education in Times of Populism and the Environmental Crisis and the Possibility of a Critical Education. In: Perspectiva 40(1):1–18. DOI: 10.5007/2175-795X.2022.e81356.
\item Kminek, H. (2023): About the need for a common and tentatively formal theory of ESD and self-critical reflections. Educational Philosophy and Theory. DOI: 10.1080/00131857.2023.2200162.
\item Singer-Brodowski, M./Kminek, H. (2023): Zu den Zielen von Bildung für nachhaltige Entwicklung und dem Stand der Implementierung im deutschen Schulsystem. DDS – Die Deutsche Schule, 115(2), 94–104. DOI: 10.31244/dds.2023.02.03.
\item Sylvester, F./Weichert, F.G./Lozano, V.L. et al. (2023): Better integration of chemical pollution research will further our understanding of biodiversity loss. Nat Ecol Evol. DOI: 10.1038/s41559-023-02117-6.
\item Kminek, H. (2023): On Education for Sustainable Development During the Lack of a Socio-ecological Transformation. (Arbeitstitel, in Vorbereitung)
\item Kminek, H./Graber, J. (2023): Engaging with Philosophy of Science in Environmental Education Research. (Arbeitstitel, in Vorbereitung)
\item Kminek, H./Schneider, F./Böhning-Gaese, K./Döll, P./Holler, H./Marg, O./Mollendorf, D./Schlottmann, A. (2023): On the Philosophy of Science Reflection of Transdisciplinary Research. A Theoretical Review. (Arbeitstitel, in Vorbereitung)
\item Kminek, H./Schink, P./Bogards, R./Böhning-Gaese, K./Curtius, J./Döll, P./Hickler, T./Holler, H./Hummel, D./Lemke, T./Lutz-Bachmann, M./Marg, O./Mollendorf, D. (2023): Planetary Thinking – A Systematic Outline of a Future Research Programme. (Arbeitstitel, in Vorbereitung)
\end{enumerate}
\subsection{Aufsätze/Beiträge in Sammelwerken und Zeitschriften [peer-reviewed]}
\begin{enumerate}
\item Kminek, H. (2015): „Plädoyer für eine auf qualitative Forschung ausgerichtete Philosophiedidaktik“. In: Zeitschrift für Didaktik der Philosophie und Ethik, 3/2015.
\item Kminek, H. (2016): „Sobre a relação entre a teoria crítica de Theodor W. Adorno e a pesquisa reconstrutiva do ensino pela hermenêutica objectiva“. In: Tagungsband: Internacional de Teoria Crítica: „Tecnologia, Violência, Memória“. Online-Publikation (2016), S. 559–567.
\item Kminek, H. (2017): „Völlig egal.“ – Bedeutungslosigkeit in einer Unterrichtsstunde zur Sprachphilosophie. In: Albus, V./Frank, M./Geier, T. (Hrsg.) (2017): Sprachliche Bildung im Philosophieunterricht. S. 162–185. Berlin/Münster.
\item Kminek, H. (2017): Prospects for a „Practise of Teaching Philosophy“. In: Oberprantacher, A./Siegetsleitner (Hrsg.) (2017): Mensch sein – Fundament, Imperativ oder Floskel? S. 713–721. Innsbruck: University Press.
\item Schindler, C./Veja, C./Kminek, H. (2017): Interfacing collaborative and multiple-layered spaces of interpretation in humanities research. In: Digital Humanities 2017, Montreal: McGill University/Université de Montréal, S. 580–583.
\item Kminek, H. (2018): „Zur grundlegenden Begründung einer empirisch gewendeten Philosophiedidaktik.“ In: Kminek, H./Thein, Chr./Torkler, R. (Hrsg.) (2018): Zwischen Präskription und Deskription. S. 13–28. Leverkusen: Budrich.
\item Pollmanns, M./Kabel, S./Leser, Chr./Kminek, H. (2018): „Krisen der Professionalisierung. Wie sich Studierende in Schulpraktischen Studien forschungsbezogenen Typs der schulischen Praxis zuwenden“. In: Artmann, M./Herzmann, P./Liegmann, A. (Hrsg.): Professionalisierung im Praxissemester.
\item Gerecht, M./Kminek, H. (2018): Veröffentlichungsfragen anderer Form: Open Data. In: Kessl, F./Jornitz, S. (Hrsg.): (Erziehung) Wissenschaftlich Publizieren. Heft 57, 02/2018. Leverkusen: Budrich.
\item Kminek, H. (2019): „Zum verfrühten, anvisierten und notwendig anzustrebendem Enden von Pädagogik“. In: Jornitz, S./Pollmanns, M. (Hrsg.) (2019): Wie mit Pädagogik enden? Leverkusen: Budrich.
\item Kabel, S./Leser, C./Pollmanns, M./Kminek, H. (2020): Fallbestimmungskrisen und Formen ihrer Bearbeitung in studentischer Fallarbeit. In: Fabel-Lamla, M./Kunze, K./Moldenhauer, A./Rabenstein, K. (Hrsg.), Kasuistik – Lehrer*innenbildung – Inklusion. Bad Heilbrunn: Klinkhardt.
\item Kminek, H. (2020): „Erziehung nach Auschwitz am Beginn des 21. Jahrhunderts. Über das Verhältnis von gesellschaftlicher Entwicklung und emanzipatorische Pädagogik“. In: Andresen, S./Nittel, D./Thompson, C. (Hrsg.) (2020): „Erziehung nach Auschwitz“.
\item Kminek, H. (2020): „Zur Lehre der Philosophiedidaktik aus der Sicht der „Philosophiedidaktik der Praxis““. In: Torkler, R. (Hrsg.) (2020): Was bedeutet: Gute Lehrerausbildung im Fach Ethik/Philosophie? Wiesbaden: Springer VS.
\item Kminek, H. (2020): „Zur „Philosophiedidaktik der Praxis“ und deren Grundlagen“. In: Thein, Chr. (Hrsg.) (2020): Philosophische Bildung und Didaktik. Wiesbaden: VS.
\item Andresen, S./Kminek, H. (2020): Erfahrungen anerkennen, Teilhabe ermoglichen, Sorgen verarbeiten. In: Pädagogik 2020 (8), S. 71–74. Weinheim: Beltz.
\item Kminek, H. (2020): „Concept of Education in Education for Sustainable Development – the Necessity of Exposing the Uncertainty“. In: Kminek/Bank/Fuchs (Hrsg.) (2020): Kontroverses Miteinander. Norderstedt: BoD.
\item Kminek, H./Wallmeier, Ph. (2020): „Nicht abschließbare Problemorientierung als Leitprinzip – Zur Bildung für die sozial-ökologische Transformation in polarisierten Zeiten“. In: Eicker, J. et al. (Hrsg.) (2020): Bildung Macht Zukunft. Wochenschau-Verlag.
\item Kminek, H. (2021): Grenzen und Möglichkeiten der Nachnutzung qualitativer Forschungsdaten aus Sicht der Objektiven Hermeneutik unter Berücksichtigung der Schul- und Unterrichtsforschung sowie der Lehrerbildung. In: Richter, C./Mojescik, K. (Hrsg.) (2021): Qualitative Sekundäranalyse. Wiesbaden: Springer VS.
\item Kminek, H. (2021): A Contribution from the Philosophy of Science for Education for Sustainable Development. In: Kminek, H./Geyer, A./Siewert, M. (Hrsg.) (2021): Engaged Reflections – Reflected Engagement. Opladen, Berlin, Toronto: Budrich.
\item Kminek, H./Holfelder, A.-K./Singer-Brodowski, M. (2021): Zukunft war gestern – Zur Legitimität der Pädagogik in Zeiten der sozial-ökologischen Krise. In: Zukunft – Stand jetzt. Jahrbuch für Pädagogik 2021. Weinheim: Beltz.
\item Kminek, H. (2022): Datenarchive in Lehre und Studium. Zur Nutzung vergessener Schatzkammern. In: Egloff, B./Richter, S. (Hrsg.): Erziehungswissenschaftlich Denken und Arbeiten. Stuttgart: Kohlhammer.
\item Kminek, H./Wahl, J. (2022): Zur (Bildungs-)Gerechtigkeit im Zusammenhang mit digitalen kollektiven Aushandlungsprozessen zur nachhaltigen Entwicklung. In: MedienPädagogik 52, S. 1–20. DOI: 10.21240/mpaed/52/2023.02.01.X.
\item Kminek, H. (2023): Education for Sustainable Development – An Aporetic Approach. ZEP – Zeitschrift für internationale Bildungsforschung und Entwicklungspädagogik, 46(2), 14–18. DOI: 10.31244/zep.2023.02.04.
\item Krug, A./Kminek, H. (2024): Bildung für nachhaltige Entwicklung im Ethik- und Philosophieunterricht – Wertediskursiv, problemorientiert und kritisch-reflexiv. In: Carrapatoso, A./Rehm, M./Reinhardt, V./Wilhelm, M. (Hrsg.), Wirksamer Unterricht in BNE (Unterrichtsqualität, Bd. 19). Baltmannsweiler: Schneider Verlag Hohengehren, S. 91–103.
\item Kminek, H. (2023): \textit{Poverty}. In: Corsaro, W. (Hrsg.): Education and Childhood Studies. Bloomsbury. (Handbook article, angenommen)
\item Kminek, H. (2023): \textit{Gender}. In: Corsaro, W. (Hrsg.): Education and Childhood Studies. Bloomsbury. (Handbook article, angenommen)
\item Dietz, T./Kminek, H./Wilmes, J.: \textit{Health, Wellbeing and Welfare}. In: Corsaro, W. (Hrsg.): Education and Childhood Studies. Bloomsbury. (Handbook article, angenommen)
\item Kminek, H./Andresen, S.: \textit{Historical Context}. In: Corsaro, W. (Hrsg.): Education and Childhood Studies. Bloomsbury. (Handbook article, angenommen)
\item Kminek, H. (2026): Zu den ethischen Fragen im Praktikum und den dialektischen Verstrickungen von Fördern und Fordern sowie von Fürsorge und Kontrolle. In: Kminek, H./Fingerle, M./Andresen, S. (Hrsg.) (2026): Erziehungswissenschaft und Ethik. Stuttgart: Kohlhammer, S. 75–92.
\item Kminek, H. (2024): Survival through Bildung in the Face of the Destruction of Human Livelihoods. In: Kminek, H. (Hrsg.): Survival through Bildung. Opladen, Berlin, Toronto: Budrich, S. 35–62.
\item Kminek, H./Bonnes, J. (2024): As Little Media Use as Possible, as Much Media Education as Necessary. On the Content-related Provisions of a Normative Demand. In: Rieckmann, M./Maurer, B./Schluchter, J. (Hrsg.): Medien – Bildung – Nachhaltige Entwicklung. Weinheim, Basel: Beltz Juventa, 02.05.2024, S. 57–68.
\item Kminek, H. (2024): Zum Beitrag der kritischen Bildungstheorie Peter Eulers im Anthropozän. In: Soziale Passagen. Berlin: Springer.
\item Kminek, H. (2025): About a Metatheory of Education for Sustainable Development as a Cross-cutting Issue of Teacher Education. In: Hermes, M./Lengler, A./Meister, N./Saß, M. (Hrsg.): Lehrkräftebildung im Wandel. Marburg: Tectum Verlag, 11.04.2025, S. 149–170.
\item Kminek, H./Holfelder, A.-K. (2025): Klimabildung im Kontext von Bildungstheorie. In: Höttecke, D./Heinicke, S./Martens, H./Nehring, A./Rabe, T. (Hrsg.): Handbuch Klimabildung. Wiesbaden: Springer VS, 20.10.2025, S. 21–44.
\item Kminek, H. (2025): Politische Bildung und BNE in situ – zu Notwendigkeit, dem Potenzial und Grenzen rekonstruktiver Unterrichtsforschung. In: Frech, S./Geyer, R./Oberle, M. (Hrsg.): Politische Bildung für nachhaltige Entwicklung. Schwalbach/Ts.: Wochenschau Verlag, 21.10.2025, S. 151–166.
\item Kminek, H. (2025): On Potentials, Limits and Concerns of Education for Sustainable Development in General, and Anthropology as Teaching Content in Particular. In: Paragrana. Internationale Zeitschrift für Historische Anthropologie. Berlin: Akademie Verlag, S. 174–185.
\item Hempel, C./Heldt, I./Kater-Wettstädt, L./Kminek, H./Bieber, G. (2026): Editorial zum Schwerpunktthema: Schulische Querschnittsaufgaben (revisited). Zur (Über-)Fachlichkeit zeitgemäßer Bildung. In: Die Deutsche Schule. Münster, New York, München, Berlin: Waxmann, S. 4–9.
\item Kminek, H. (2026): About the Further Development of a Metatheoretical Analysis Grid for the Field of Education and (Non-)Sustainability. In: Education Sciences. Basel: MDPI Publishing, S. 1–23.
\item Ehras, C./Singer-Brodowski, M./Costa, J./Hemmer, I./Brenne, A./Denzin, A./Dick, M./Fischer, D./Kenner, S./Kminek, H./Kreid, S./Krengel, F./Lindau, A./Lohmann, J./Menthe, J./Meyer, M./Roczen, N./Schaal, S./Surkamp, C./Tacke, L./Weselek, J. (2026): Professionalisation of Teachers for Education for Sustainable Development in Germany: A Scoping Review. In: RISTAL – Research in Subject-matter Teaching and Learning, S. 1–64.
\end{enumerate}
\subsection{Rezensionen/Berichte/Vermischtes}
\begin{enumerate}
\item „Twardella, Johannes (2015): Pädagogische Kasuistik: Fallstudien zu grundlegenden Fragen des Unterrichts. Opladen/Berlin/Toronto: Budrich. 227 Seiten.“ In: Zeitschrift für Didaktik der Philosophie und Ethik, 38 (4), S. 115–116.
\item „Samuel Scheffler (2015): Der Tod und das Leben danach. Berlin: Suhrkamp. 155 Seiten.“ In: Zeitschrift für Didaktik der Philosophie und Ethik, 39 (1), S. 101.
\item „Gold, Andreas (2015): Guter Unterricht: Was wir wirklich darüber wissen. Göttingen: Vandenhoeck \& Ruprecht. 176 Seiten.“ In: Zeitschrift für Didaktik der Philosophie und Ethik, 39 (1), S. 95–96.
\item „Geier, Thomas/Pollmanns, Marion (Hrsg.) (2016): Was ist Unterricht? Zur Konstitution einer pädagogischen Form. Wiesbaden: Springer. 252 Seiten.“ In: Zeitschrift für Didaktik der Philosophie und Ethik, 39 (1), S. 95–96.
\item Projektbericht CEDIFOR (2017). Abschnitt: Forschungskapazitäten für die qualitative Forschung durch Kollaboration und semantische Auszeichnung. Das Beispiel Unterrichtsinteraktion.
\item „Klaus Feldmann (2017): Handelndes Lernen im Philosophieunterricht. Wiesbaden: Springer VS. 212 Seiten.“ In: Zeitschrift für Didaktik der Philosophie und Ethik.
\item Abschlussbericht Förderfonds Lehre (2018): Unterstützung von studienzentrierter forschender Lehre mittels einer VFU für Objektive Hermeneutik.
\item Kann Pädagogik Waldbrände verhindern? Zwischen Zurichtung und Mündigkeit: Bildung im Sinne einer nachhaltigen Entwicklung. In: FORSCHUNG FRANKFURT. Das Wissenschaftsmagazin der Goethe-Universität, 2/2020, S. 74–78.
\item Kminek, H./Singer-Brodowski, M./Costa, J./Weselek, J./Ruckelshauß, T. (2024): Tagungsbericht: Kontroversen und Debatten im Kontext von Bildung für nachhaltige Entwicklung. In: ZEP – Zeitschrift für internationale Bildungsforschung und Entwicklungspädagogik. Münster, New York, München, Berlin: Waxmann, S. 30.
\end{enumerate}
\section{Vorträge}
{\itshape\small\color{cvgrey}(Die mit „*“ gekennzeichneten Beiträge durchliefen ein double-blind Peer-Review-Verfahren.)\par}\vspace{4pt}
\begin{enumerate}
\item „Wie zwischen Überhöhung und Abwertung Philosophie zerrieben wird, ohne sie aufzulösen“. 05.04.2014. Kolloquium Philosophiedidaktik. Frankfurt am Main.
\item „Was geschieht mit den Daten bei ApaeK?“. 10.04.2014. Deutsche Gesellschaft f. Soziologie, Sektion Biografieforschung. Frankfurt am Main.
\item (Philosophische) Anthropologie im und in Unterricht. 07.10.2014. Menschenbilder in Schule und Unterricht. Trier.
\item „Empirische Erschließungen von Philosophie-Unterricht an Schulen im Anschluss an Adorno für Philosophie“. 05.12.2014. Doktorandenkolloquium der Österreichischen Gesellschaft für Philosophie. Innsbruck (Österreich).
\item „Unausgewertete Ergebnissicherung […]“. 15.01.2015. Unterrichtsergebnisse: Erziehungswissenschaftliche und fachdidaktische Perspektiven. Mainz.
\item „Methodisch verhindertes Philosophieren – eine Fallanalyse“. 31.03.2015. Arbeitstagung Doktorand\_innen und Habilitand\_innen der Philosophiedidaktik. Frankfurt am Main.
\item „Realistische Fachdidaktik am Beispiel des Philosophie-Unterrichts“. 09.04.2015. Kolloquium der Fachdidaktiken. Frankfurt am Main.
\item „Empirisch-qualitative Erschließung von Philosophie-Unterricht an Schulen als Aufklärung der und für die Praxis“. 04.–06.06.2015. Kongress der Österreichischen Gesellschaft für Philosophie. Innsbruck (Österreich).
\item „Philosophiefachdidaktik der Praxis – Skizze eines neuen Forschungsparadigmas“. 01.07.2015. Öffentliche Abendvorträge zu Philosophiedidaktik und Bildungsphilosophie. Mainz.
\item „Wie wird Kants Text ‘Beantwortung der Frage: Was ist Aufklärung?‘ im Unterricht verhandelt?“. 13.07.2015. Fachverband Philosophie e.V. Hessen. Frankfurt am Main.
\item „Analysen studentischer Fallarbeit: Zur Wirklichkeit der Professionalisierung in Schulpraktischen Studien forschungsbezogenen Typs.“ 22.04.2016. 4. Tagung der Arbeitsgemeinschaft Kasuistik in der Lehrer\_innenbildung. Göttingen. (gemeinsam mit Prof.’in Dr. M. Pollmanns, Dr. Chr. Leser, S. Kabel und R. Hünig)
\item „Education, Politics and Negative Dialectics“. 02.–04.06.2016. Critical Theory Conference: Adorno and Politics. Istanbul (Türkei).
\item „Zum Verhältnis der kritischen Theorie Theodor W. Adornos zur rekonstruktiver Unterrichtsforschung mittels Objektiver Hermeneutik“. 10.–14.10.2016. Tecnologia, Violência, Memória. São Carlos (Brasilien).
\item „Krisen studentischer Fallarbeit und Formen ihrer Bewältigung […]“. 06.–08.03.2017. 2. Internationaler Kongress „Lernen in der Praxis“ der IGSP. Bochum. (gemeinsam mit Prof.’in Dr. M. Pollmanns, Dr. Chr. Leser und S. Kabel)
\item „Heydorn’s Critical Theory of Education“. 18.–20.10.2017. IV International Congress of Critical Theory. Madrid (Spanien).
\item *„We all have to protect the environment, but somehow nobody does it.“ 07.–09.11.2017. Colloque «Changements et Transitions: enjeux pour les éducations à l’environnement et au développement durable». Toulouse (Frankreich).
\item *„Education for maturity and responsibility in the beginning of the 21st century“. 10.–12.05.2018. 11th International Critical Theory Conference of Rome. Rome (Italy).
\item „Education for maturity and responsibility in the beginning of the 21st century“. 17.–19.05.2018. Modernity between Damaged Life and Sane Society. Chicago (USA).
\item „Philosophie und Philosophieren im Unterricht. Ein Beitrag aus der empirischen Unterrichtsforschung.“ 29.06.–01.07.2018. 3. Internationale Arbeitstagung für Didaktik der Philosophie und Ethik. Universität zu Köln.
\item „About Education in Times of Populism and the Possibility of a Critical Education“. 23.–25.01.2019. Fascism? Populism? Democracy? University of Brighton (UK).
\item *„Zur Bestimmung der Grenzen der politischen Bildung.“ 27. DGfE-Kongress OPTIMIERUNG, 15.–18.03.2020, Köln (pandemiebedingt ausgefallen).
\item *„About Education for Sustainable Development and its Aporias.“ ECER 2020, Glasgow, 25.–28.08.2020 (pandemiebedingt ausgefallen).
\item Bildung für nachhaltige Entwicklung als Forschungsgegenstand und -praxis der Grenzbearbeitung. Kommentar. Jahrestagung der Sektion SIIVE in der DGfE, 20.02.2021.
\item Bildung für eine nachhaltige Entwicklung – eine aporetische Operation. Eingeladener Online-Vortrag, Flensburg, 15.06.2021.
\item Symposium „Analyse und Reflexion von Grenzbearbeitungen in erziehungswissenschaftlichen Forschungsprozessen“, DGfE-Kongress 2022, Kommentar.
\item Bildung und Nachhaltigkeit. PH Heidelberg, 26.06.2023 (eingeladener Vortrag).
\item *Education for Sustainable Development: Learning From, Through and About Social Movements for Climate Action, EERA 2023, Glasgow (Chair and discussant).
\item Fachsitzung: Planetary Thinking for Planetary Futures. Deutscher Kongress für Geographie, 21.09.2023, Frankfurt am Main (Sitzungsleitung gemeinsam mit Prof.’in Dr. P. Döll und eigener Beitrag).
\item *„Jetzt sind aber wir nun mal […] in nem Wirtschaftssystem gefangen, was eben nicht nach Postwachstumsgedanken agiert“ – zu wissenschaftstheoretischen und forschungsmethodologischen Fragen im Konnex von Nachhaltigkeit und Erwachsenenbildung. Jahrestagung der DGfE-Sektion Erwachsenenbildung, 11.–13.09., München.
\end{enumerate}
\cvausblick{Ausblick (Stand 2023)}
\begin{enumerate}
\item *On the research questions of an ESD in view of the latest socioecological diagnoses and research results. DGfE-BNE-Kommissionstagung, 25.–27.09.2023, PH Heidelberg.
\item *Symposium: „Das Anthropozän als planetarische Krise: Global Citizenship Education und die Transformation unseres Verständnisses vom Menschen“. DGfE-Kongress 2024, Halle/Saale (Diskutant; angenommen).
\item *Working Group: Sustainability, Crisis and Transformation. DGfE-Kongress 2024, Halle/Saale (Chair, contribution and discussant; reserve list).
\end{enumerate}
\cvsubhead{Vorträge 2024–2026 (Alpen-Adria-Universität Klagenfurt)}
\begin{enumerate}
\item Possibilities and Limitations of ESD Research. 23.05.2024. Blended Intensive Programme ‘Environmental and Sustainability Education‘, Universität Vechta. Vortrag auf Einladung.
\item Zur Bedeutung von Kunst und Kultur in der Bildung für nachhaltige Entwicklung. 27.05.2024. Green Talks, Frankfurt am Main. Vortrag auf Einladung.
\item About the Potential and Limits of Education for Sustainable Development. 12.09.2024. New Frontiers in Climate Ethics and Politics, Mailand (Italien). Vortrag auf Einladung.
\item Zum Nexus von Bildung für nachhaltige Entwicklung und Macht und dessen Systematisierung. 18.09.2024. „Spannungsfelder von Bildung für nachhaltige Entwicklung und die zunehmende Ausdifferenzierung des Feldes“, Bamberg. Angemeldeter Vortrag.
\item Zum politischen Anspruch eines transformativen BNE-Dialogs zur Nachhaltigkeitsbildung als postpolitische Praxis und mögliche Orientierungen für die Repolitisierung von BNE im Allgemeinen und dem GW-Unterricht im Besonderen. 07.04.2025. Zukunft Fachdidaktik GW – 2025 Transformative politische Bildung in GW in Zeiten multipler Krisen, Schlierbach. Vortrag auf Einladung.
\item About the Nexus of Emotions and Elements of a Theory Architecture for Environmental and Sustainability Education (ESE). 11.04.2025. EMESER Symposium 2025 (online). Angemeldeter Vortrag.
\item Zu den Grenzen und Möglichkeiten von Bildung für nachhaltige Entwicklung. 21.05.2025. Lectures for Future, Universität für Bodenkultur Wien. Vortrag auf Einladung.
\item Zur Geschichte und ausgewählter Modelle Transdisziplinärer Forschung. 23.05.2025. Forschungsworkshop „Wochenende der Forschungsfragen“ (Kooperation SCiENCE\_LINKnockberge 2025), UNESCO Biosphärenpark Kärntner Nockberge. Vortrag auf Einladung.
\item „Bewußtsein ist alles“ – Zum Menschenbild in der Bildungsphilosophie Heinz Joachim Heydorns. 12.07.2025. „Pädagogisch-anthropologische Grenzziehungen: Menschenbilder in der Bildungswissenschaft“, Innsbruck. Angemeldeter Vortrag.
\item About Systematic (Meta-)Theory Formation as an Answer to the Question of ESER’s Identity. 09.09.2025. ECER 2025, Belgrad. Angemeldeter Vortrag.
\item About Criticising of, within and for ESER Discourse. 11.09.2025. ECER 2025, Belgrad. Angemeldeter Vortrag.
\item „We all have to protect the environment, but somehow nobody does.“ – Or: About the Potential and Limits of Education for Sustainable Development. 22.09.2025. Course Sustainability, CIEE Berlin. Vortrag auf Einladung.
\item Kritische Reflexionen zum Forschungsprogramm einer Bildung für nachhaltige Entwicklung im Nexus zu Demokratiepädagogik und Politischer Bildung. 25.09.2025. Bildung, Demokratie \& Nachhaltigkeit – Jahrestagung 2025 der DGfE-Kommission Bildung für nachhaltige Entwicklung, Dresden. Angemeldeter Vortrag.
\item Zu den Aufgaben der AG „Nachhaltigkeit bzw. Nicht-Nachhaltigkeit und planetare Zukünfte“ – kritische Reflexionen und Impulse. 30.09.2025. Auftakttagung der DGfE-Arbeitsgemeinschaft „Nachhaltigkeit bzw. Nicht-Nachhaltigkeit und planetare Zukünfte“, Köln. Angemeldeter Vortrag.
\item About the ‘Sustainable Nonsustainability‘ – Asking Questions as Part of the Answers. 06.11.2025. Colloque Enjeux de la coopération internationale environnementale, Université Catholique de Lille (Frankreich). Angemeldeter Vortrag.
\item Zur Frage nach einer angemessenen Grundbildung in der Dimension des Wissens im Anthropozän. 23.03.2026. DGfE-Kongress 2026 „Brüche“, München. Angemeldeter Vortrag.
\item Kontinuitäten – Brüche – Vereinigungen. Reflexionen zur erziehungswissenschaftlichen Theoriebildung am Beispiel Bildung für nachhaltige Entwicklung. 25.03.2026. DGfE-Kongress 2026 „Brüche“, München. Vortrag auf Einladung.
\item Educational science: Education for sustainable development – (Meta-)theoretical reflections on the state of research and research desiderata. 20.04.2026. Lectures for Future II, TU Wien. Vortrag auf Einladung.
\item Zum Beitrag Schleiermachers für eine Metatheorie und ein Forschungsprogramm der Erziehungswissenschaft. 30.04.2026. Schleiermacher-Tagung 2026: Was uns zusammenhält – Über Möglichkeiten und Grenzen der Erziehung als Grundlage sozialen Zusammenhalts, Linz. Angemeldeter Vortrag.
\item On the relationship between general and specific subject didactics regarding education for sustainable development. 20.05.2026. Didactica Viva: General and Subject-specific Didactics as Living and Responsive Fields, Brno (Tschechien). Angemeldeter Vortrag.
\item Expertise – Emancipation – Maturity. On the question of contemporary political education. 22.05.2026. Political Education as Emancipation, Wien. Angemeldeter Vortrag.
\item How objective is, can, and should ESD research be? 03.06.2026. Education for Sustainable Development – Controversial?! Systematizing the Research Field, Universität Regensburg (digital). Vortrag auf Einladung.
\end{enumerate}
\subsection{Poster}
\begin{enumerate}
\item Virtuelle Forschungsumgebung – eine Bereicherung für Forschung und Lehre? 03.–04.11.2017. Fachtagung der DGfE: „Universität 4.0“. Berlin. (gemeinsam mit Dr. C. Schindler und Dr. M. Meier)
\end{enumerate}
\subsection{Workshops}
\begin{enumerate}
\item „Möglichkeiten und Grenzen der Verbindung von Fachdidaktik und Kasuistik […]“ (gemeinsam mit JProf. Dr. A. Burkard und M. Frank). 3.–4.03.2017. 6. Arbeitstagung AG Kasuistik in der Lehrer\_innenbildung. Halle (Saale).
\item „Virtuelle Forschungsumgebung – eine Bereicherung für Forschung und Lehre?“. 22.–23.04.2016. 4. Tagung AG Kasuistik in der Lehrer\_innenbildung. Göttingen.
\item „Digitalisierte Hochschullehre zu qualitativen Forschungsmethoden – zum Beispiel Semantic Cora […]“. 28.–29.06.2018. Universität Duisburg-Essen.
\item „Ethische Fragen im Zusammenhang mit der Praktikumsbetreuung“. Tag der Forschungsethik am Fachbereich Erziehungswissenschaften, 07.05.2019, Universität Frankfurt.
\end{enumerate}
\section{Veranstaltungen – Organisation und Mitwirkung}
\begin{longtable}{|>{\small\raggedright\arraybackslash}p{3.65cm}|>{\small\raggedright\arraybackslash}p{4.76cm}|>{\small\raggedright\arraybackslash}p{2.22cm}|>{\small\raggedright\arraybackslash}p{4.84cm}|}
\hline
\rowcolor{navy}
\textcolor{white}{\textbf{Veranstaltungsart}} & \textcolor{white}{\textbf{Thema}} & \textcolor{white}{\textbf{Zeit/Ort}} & \textcolor{white}{\textbf{Tätigkeit}} \\
\hline
\endfirsthead
\hline
\rowcolor{navy}
\textcolor{white}{\textbf{Veranstaltungsart}} & \textcolor{white}{\textbf{Thema}} & \textcolor{white}{\textbf{Zeit/Ort}} & \textcolor{white}{\textbf{Tätigkeit}} \\
\hline
\endhead
\hline
\endfoot
Kolloquium & Kolloquium der Fachdidaktiken der Goethe-Universität Frankfurt am Main (halbjährlich) & 2013–2015, Frankfurt & Organisation, Durchführung, eigener Vortrag \\
\hline
\rowcolor{lightgrey}
Arbeitstagung & 1.–5. Arbeitstagung für Doktorand\_innen und Habilitand\_innen der Philosophiedidaktik (jährlich) & 2014–2018, Frankfurt & Organisation, Durchführung, eigener Vortrag \\
\hline
Int. Arbeitstagung Philosophiedidaktik & 1. Internationale Arbeitstagung der Philosophiedidaktik: Zwischen Präskription und Deskription & Juni 2016 & Organisation, Durchführung, eigener Vortrag \\
\hline
\rowcolor{lightgrey}
Kommissionstagung Grundschulforschung (DGfE) & Grundschulforschungstagung 2018: Diversität und soziale Ungleichheit & Sept. 2018 & Gutachtertätigkeit \\
\hline
Kommissionstagung Schulforschung u. Didaktik (DGfE) & Transnationale Perspektiven auf Schule und Bildung & Sept. 2018 & Mitglied des Organisationsteams, inkl. Begutachtung der Einreichungen \\
\hline
\rowcolor{lightgrey}
Arbeitstagung & Auftakttagung Zeitschriftenneugründung: Datum \& Diskurs & März 2019, EUF Flensburg & Mitglied des Organisationsteams \\
\hline
Arbeitstagung & 11. Tagung der AG Kasuistik in der Lehrer\_innenbildung (SIG der DGfE) & Herbst 2019, Universität Mainz & Mitglied des Organisationsteams, inkl. Begutachtung der Einreichungen \\
\hline
\rowcolor{lightgrey}
Vortragsreihe & Planetary Thinking (interdisziplinäre Arbeitsgruppe, Goethe-Universität) & 2021–2022 & Mitglied des Organisationsteams, eigener Vortrag (02/2022) \\
\hline
Workshop & Planetary Boundaries – between Facts and Norms (mit Forschungskolleg Humanwissenschaften, Bad Homburg) & 07/2021 & Mitglied des Organisationsteams \\
\hline
\rowcolor{lightgrey}
DGfE-BNE-Kommissionstagung 2022 & „Ökologische, gesellschaftliche und individuelle Umbrüche und ihre Bedeutung für BNE“ & Sept. 2022 & Organisation, Durchführung \\
\hline
DGfE-BNE-Kommissionstagung 2023 & „Kontroversen und Debatten im Kontext von Bildung für nachhaltige Entwicklung“ & Sept. 2023 & Mitglied des Organisationsteams, inkl. Begutachtung der Einreichungen \\
\hline
\rowcolor{lightgrey}
Tagung/Konferenz & Areas of tension between education for sustainable development and the increasing differentiation of the field & 18.–20.09.2024, Bamberg & Mitglied des Organisationsteams (mit J. Costa, A.-K. Holfelder, J. Hufnagl, M. Osterrieder, A. Scheunpflug, M. Singer-Brodowski, D. Taube) \\
\hline
Workshop & DGfE Dialogue Workshop: „Sustainability and planetary futures as a transformation challenge for educational science“ & 14.–15.11.2024, Nürnberg & Mitglied des Organisationsteams (mit J. Costa, M. Singer-Brodowski, Ch. Wulf, U. Stenger, L. Klepacki, N. Grünberger, J. Engel) \\
\hline
\rowcolor{lightgrey}
Tagung/Konferenz & Transnationalization of social and ecological inequalities and their relevance to educational science & 20.–21.03.2025, Berlin (TU) & Mitglied des Organisationsteams (mit D. Amirpur, N. Bernhard, J. Costa, A. Dogmus, K. Falkenberg, S. Hartong, R. Nikolai, J. Panagiotopoulou, S. Ress, M. Singer-Brodowski) \\
\hline
DGfE-BNE-Kommissionstagung 2025 & Education, Democracy \& Sustainability – Jahrestagung 2025 der DGfE-Kommission Bildung für nachhaltige Entwicklung & 23.–25.09.2025, Dresden & Mitglied des Organisationsteams (mit J. Costa, A.-K. Holfelder, J. Markert, A. Goller) \\
\hline
\rowcolor{lightgrey}
Tagung/Konferenz & Erziehungswissenschaft und Nachhaltigkeit – Widerstände bearbeiten. Zukünfte gestalten. Auftakttagung der DGfE-AG „Nachhaltigkeit bzw. Nicht-Nachhaltigkeit und planetare Zukünfte“ & 29.09.–01.10.2025, Köln & Mitglied des Organisationsteams (mit M. Singer-Brodowski, U. Stenger, Ch. Wulf, O. Bilgi, J. Dinkelaker, J. Elven, S. Gaubitz, L. Klepacki, T. Klepacki, S. Ress, O. Sanders) \\
\hline
\end{longtable}
\section{Funktionen und Gutachtertätigkeiten}
\subsection{Editorial Advisor}
\cvitem{2018}{Journal of Decolonising Disciplines.}

\subsection{Mitglied des Editorial Board}
\cvitem{2016}{Journal of Didactics of Philosophy.}
\cvitem{seit 11/2025}{Educations (openscience.fr/Educations).}

\subsection{Editorial-/Herausgebertätigkeit bei Zeitschriften und Schriftenreihen}
\cvitem{seit 04/2024}{Editor – datum \& diskurs – Zeitschrift für Komplementär-, Sekundär- und Reanalysen im Feld der qualitativen Schul- und Unterrichtsforschung.}
\cvitem{seit 04/2024}{Editor – Wissenschaftliche Beiträge zur Philosophiedidaktik und Bildungsphilosophie (Scientific Contributions to Philosophy Didactics and Educational Philosophy).}

\subsection{Gutachtertätigkeiten für Zeitschriften [double-blind peer-reviewed]}
\begin{itemize}
\item Childhood Vulnerability Journal
\item Diskurs Kindheits- und Jugendforschung / Discourse. Journal of Childhood and Adolescence Research
\item Environmental Education Research
\item Herausforderung Lehrer*innenbildung – Zeitschrift zur Konzeption, Gestaltung und Diskussion
\item Journal of Didactics of Philosophy
\item Qualifizierung für Inklusion. Zeitschrift zur Forschung über Aus-, Fort- und Weiterbildung pädagogischer Fachkräfte
\item Zeitschrift für Medienpädagogik
\item Zeitschrift für Pädagogik
\end{itemize}
\subsection{Weitere Gutachtertätigkeiten für Zeitschriften}
\cvitem{04–04/2024}{The Journal of Environmental Education.}
\cvitem{11–12/2024}{Journal for Educational Science.}
\cvitem{04/2025}{Discover Education.}
\cvitem{11/2025}{ZEP – Zeitschrift für internationale Bildungsforschung und Entwicklungspädagogik.}
\cvitem{11–12/2025}{Foundations of Science.}
\cvitem{01/2026}{Discover Sustainability.}
\cvitem{02–03/2026}{Discover Sustainability.}
\cvitem{02–05/2026}{Journal of Curriculum Studies.}
\cvitem{03/2026}{Journal of Didactics of Philosophy.}
\cvitem{03/2026}{Ludwigsburg Contributions to Media Education (LBzM).}
\cvitem{05/2026}{Education Sciences.}
\cvitem{05/2026}{Sustainability.}
\cvitem{05/2026–04/2027}{Sustainability, Education Sciences.}
\cvitem{05–06/2026}{Frontiers in Sociology.}
\cvitem{06/2026}{Sustainability.}
\cvitem{06/2026}{Revue Education.}

\subsection{Gutachtertätigkeiten für Sammelwerke und Handbücher}
\cvitem{02–04/2025}{Nachhaltige Transformation der Wirtschaft (Hrsg.: B. Weber/S. von Lingen/H. Hantke).}
\cvitem{02–04/2025}{Handbuch Bildung für nachhaltige Entwicklung (Hrsg.: F. Bertschy/M. Rieckmann/M. Singer-Brodowski).}

\subsection{Funktionen bei Konferenzen und Tagungen}
\cvitem{05–07/2025}{Reviewer – DGfE-Kongress 2026 „Brüche“.}
\cvitem{09/2025}{Moderation – ECER 2025.}
\cvitem{12/2025}{Reviewer – DGfE-Kongress 2026 „Brüche“.}
\cvitem{01/2026}{Reviewer – ECER.}
\cvitem{03/2026}{Reviewer – ÖFEB-Kongress 2026 „Education in Times of Glocalization“.}

\subsection{Funktionen in Gremien und Arbeitsgruppen}
\cvitem{04/2024–03/2025}{Vorsitzender (Chair) der DGfE-Kommission Bildung für nachhaltige Entwicklung.}
\cvitem{04/2024–03/2027}{Mitglied des Vorstands der Sektion Interkulturelle und International Vergleichende Erziehungswissenschaft (DGfE).}
\cvitem{04/2024–11/2025}{Mitglied des Scientific Advisory Board, Projekt ENCOMPASS.}
\cvitem{seit 06/2025}{Mitglied der AG „Gemeinsames Masterstudium der AAU und Universität Leoben zum Thema ‘Circular Engineering‘ / ‘Circular Business‘“.}
\cvitem{seit 04/2026}{Reviewer, Arbeitsgruppe Wissenschaftliche Integrität (AAU).}

\subsection{Weitere Gutachtertätigkeit}
\begin{itemize}
\item APCC Special Reports: Strukturen für ein klimafreundliches Leben! (sr22.ccca.ac.at)
\item Zahlreiche Tagungen und Konferenzen (u. a. ECER 2023; weitere ab 2024 s. oben)
\end{itemize}
\section{Eingeworbene Dritt- und Fördermittel}
{\itshape\small\color{cvgrey}(* = Mitantragsteller, sonst Hauptantragsteller)\par}\vspace{4pt}
\begin{longtable}{|>{\small\raggedright\arraybackslash}p{2.22cm}|>{\small\raggedright\arraybackslash}p{5.87cm}|>{\small\raggedright\arraybackslash}p{4.29cm}|>{\small\raggedright\arraybackslash}p{3.09cm}|}
\hline
\rowcolor{navy}
\textcolor{white}{\textbf{Zeitraum}} & \textcolor{white}{\textbf{Thema}} & \textcolor{white}{\textbf{(Dritt-)Mittelgeber}} & \textcolor{white}{\textbf{Betrag}} \\
\hline
\endfirsthead
\hline
\rowcolor{navy}
\textcolor{white}{\textbf{Zeitraum}} & \textcolor{white}{\textbf{Thema}} & \textcolor{white}{\textbf{(Dritt-)Mittelgeber}} & \textcolor{white}{\textbf{Betrag}} \\
\hline
\endhead
\hline
\endfoot
* 2016–2017 & Pilotstudie: Forschungskapazitäten für die qualitative Forschung durch Kollaboration und semantische Auszeichnung & BMBF & 30.000 € \\
\hline
\rowcolor{lightgrey}
2017 & QSL-Förderfonds Lehre – VFU für Obj. Hermeneutik & Universität Frankfurt & 1.680 € \\
\hline
2018 & Projekt aKodi – Studierenden- und handlungsorientierte Lehre im Kontext der Digitalisierung & Universität Frankfurt & 1.680 € \\
\hline
\rowcolor{lightgrey}
2018 & QSL-Förderfonds Lehre – Ringvorlesung „Bildung für nachhaltige Entwicklung“ & Universität Frankfurt & 4.930 € \\
\hline
2019 & Bildung für nachhaltige Entwicklung durch die und an der Goethe-Universität & Universität Frankfurt & 15.670 € \\
\hline
\rowcolor{lightgrey}
2020 & „Konjunkturpaket als Motor für eine sozial-ökologische Transformation?“ & Grade Sustain, Universität Frankfurt & 2.844 € \\
\hline
2020 & „Empowering for the Transition Towards a Climate Neutral, Sustainable Europe by Reasons“ & Förderfonds Aufbau koordinierter Programme, Universität Frankfurt & 4.996 € \\
\hline
\rowcolor{lightgrey}
* 2021 & Learning Democracy in Europe & BMBF & 44.923,66 € \\
\hline
2021 & „Sustainable Europe by Education of Deliberation“ (SEED) & European Union & 4.851.498,38 € \\
\hline
\rowcolor{lightgrey}
2022 & DGfE-BNE-Kommissionstagung 2022 & Goethe-Universität GRADE SUSTAIN & 5.000,00 € \\
\hline
2022 & DGfE-BNE-Kommissionstagung 2022 & FB Erziehungswissenschaften, Goethe-Universität & 1.500,00 € \\
\hline
\rowcolor{lightgrey}
2022 & Transformationspfade und wissenschaftstheoretische Reflexionen & RobustNature SynergyFund, Goethe-Universität & 26.764 € \\
\hline
2023 & Goethe Teaches Sustainability & Goethe-Universität & 146.511,84 € \\
\hline
\rowcolor{lightgrey}
2023 & Zur wissenschaftstheoretisch-reflektierten Systematisierung von BNE & DFG & ca. 50.000 € (in Vorbereitung) \\
\hline
\end{longtable}
\section{Gremienarbeit}
\cvitem{seit 03/2021}{Vorstandsmitglied der Kommission Bildung für eine nachhaltige Entwicklung in der Deutschen Gesellschaft für Erziehungswissenschaft.}
\cvitem{seit 2018}{Mitglied der Arbeitsgruppe „Internationalisierung“ am Fachbereich Erziehungswissenschaften der Goethe-Universität Frankfurt am Main.}
\cvitem{2020–2023}{Mitglied der vom Senat der Goethe-Universität eingesetzten Arbeitsgruppe „Nachhaltigkeit in und an der Goethe-Universität“.}
\cvitem{seit 2016}{Stellvertretender Vorsitzender der Ethikkommission des Fachbereichs Erziehungswissenschaften der Goethe-Universität Frankfurt am Main (04/2020–03/2021 kommissarische Leitung).}
\cvitem{2017–2019}{Gewähltes Mitglied der Mittelbauvertretung am Fachbereich Erziehungswissenschaften der Goethe-Universität Frankfurt am Main.}
\cvitem{2016–2017}{Mitarbeiter in einer Berufungskommission an der Goethe-Universität Frankfurt.}
\cvitem{2010–2012}{Mitglied des Seminarrates am Studienseminar Offenbach.}
\cvitem{2007–2009}{Mitglied des L-Netzes der Universität Frankfurt.}
\cvitem{2003–2005}{Mitglied der Fachschaft Erziehungswissenschaften der Universität Frankfurt.}

\section{Fellowships}
\cvitem{2014}{Doktorandenkolloquium Österreichische Gesellschaft für Philosophie, Innsbruck (Österreich).}
\cvitem{2016}{Vortragsreise nach Brasilien (DAAD gefördert).}
\cvitem{2018}{Austauschdozent PH Bern (Swiss European Mobility Programme), Schweiz, im November 2018.}
\cvitem{2019}{Austauschdozent Alpen-Adria-Universität Klagenfurt (ERASMUS+-Programm), Österreich, im März 2019.}

\section{Mitgliedschaften und außeruniversitäres Engagement}
\begin{itemize}
\item Mitglied der Deutschen Gesellschaft für Erziehungswissenschaften – Sektion Schulpädagogik; AG Kasuistik in der Lehrerbildung; Interkulturelle und International Vergleichende Erziehungswissenschaft; Kommission Bildung für nachhaltige Entwicklung (seit 02/2021 Mitglied des Vorstands).
\item Gesellschaft für Philosophie- und Ethikdidaktik (Gründungsmitglied).
\item Deutsche Gesellschaft für Philosophie.
\item Vorsitzender / Stellvertretender Vorsitzender / Schriftführer, Fachverband Philosophie e. V. Hessen (2014–2022).
\item Mitglied der Gesellschaft für Bildung und Wissen (seit 2010).
\item Mitglied der Gewerkschaft Erziehung und Wissenschaft (GEW) (2010–2022).
\item Mitglied der Vereinten Dienstleistungsgewerkschaft (ver.di) (2006–2010).
\end{itemize}

